% !TeX root = part4.tex
\documentclass[../final.tex]{subfiles}

\begin{document}
\ifSubfilesClassLoaded{\setcounter{chapter}{3}\setcounter{section}{7}}{}

% Лекция 22.09. Источники: фотографии 1.jpg--16.jpg.
% Проверка правил E1/M1/E2: https://physics.nist.gov/Pubs/AtSpec/node17.html
% Все схемы построены в TikZ; масштаб энергии условный, кроме подписанных чисел.

\section{Правила отбора для многоэлектронного атома}

\rconcept{Правила отбора} определяются матричным элементом оператора
перехода:
\begin{equation*}
    V_{fi}=\langle f|\hat V|i\rangle,\qquad
    w_{fi}\propto|V_{fi}|^2,\qquad V_{fi}=0\ \Rightarrow\ w_{fi}=0.
\end{equation*}

\subsection{Одноэлектронный характер дипольного перехода}

При неподвижном ядре дипольный момент электронной оболочки равен
\begin{equation*}
    \hat{\mathbf d}=q_e\sum_a\hat{\mathbf r}_a,\qquad q_e=-e,\quad e>0.
\end{equation*}
Для двух электронов в ортонормированных орбиталях:
\begin{align*}
    \langle a'b'|\hat{\mathbf d}|ab\rangle
    =q_e\bigl[&\langle a'|\hat{\mathbf r}|a\rangle\delta_{b'b}
    +\delta_{a'a}\langle b'|\hat{\mathbf r}|b\rangle\bigr],\\
    a'\ne a,\quad b'\ne b
    &\quad\Longrightarrow\quad \langle a'b'|\hat{\mathbf d}|ab\rangle=0.
\end{align*}
Следовательно, \rresult{в E1-переходе меняется не более одной занятой
спин-орбитали}. Это верно и для детерминантов Слейтера.
Условие применимости --- \rassumption{пренебрежение смешиванием конфигураций}.

\subsection{Мультипольность и строгие правила}

E1, M1 --- электрический и магнитный диполь; E2, M2 --- квадруполи.
Ранг оператора \(k=1\) для диполя и \(k=2\) для квадруполя:
\begin{equation*}
    \mathrm{E1}:\ \mathbf d,\qquad
    \mathrm{M1}:\ \boldsymbol\mu,\qquad
    \mathrm{E2}:\ Q_{uv}\sim\sum_a q_e r_{au}r_{av},\qquad
    \mathrm{M2}:\ M_{uv}\sim\sum_a\mu_{au}r_{av}.
\end{equation*}
Для квадруполей подразумеваются тензорные части ранга 2;
\(a\) нумерует электроны, \(u,v\) --- координаты.

Для атома без внешнего поля:
\begin{equation*}
    |J_f-J_i|\le k\le J_f+J_i,\qquad
    |\Delta M_J|\le k,\qquad
    \pi=(-1)^{\sum_a\ell_a}.
\end{equation*}
Здесь \(\pi\) --- чётность конфигурации, первое условие --- правило
треугольника для моментов.

\begin{center}
\small\renewcommand{\arraystretch}{1.4}
\begin{tabular}{l|cc|cc}
    \hline
    \rconcept{Правило}&\rconcept{E1}&\rconcept{M1}&\rconcept{E2}&\rconcept{M2}\\
    \hline
    \(\Delta J\)&\multicolumn{2}{c|}{\(0,\pm1\)}&\multicolumn{2}{c}{\(0,\pm1,\pm2\)}\\
    Исключения&\multicolumn{2}{c|}{\(0\not\leftrightarrow0\)}&
    \multicolumn{2}{c}{\(0\not\leftrightarrow0,1;\quad\frac12\not\leftrightarrow\frac12\)}\\
    \(\Delta M_J\)&\multicolumn{2}{c|}{\(0,\pm1\)}&\multicolumn{2}{c}{\(0,\pm1,\pm2\)}\\
    Чётность&меняется&сохраняется&сохраняется&меняется\\
    \hline
\end{tabular}
\end{center}
Для отдельных компонент учитывается поляризация.
При дипольном переходе с \(\Delta J=0\) компонент \(M_i=M_f=0\) отсутствует.

\subsection{Конфигурация и LS-связь}

Без смешивания конфигураций и при \(\Delta S=0\):
\begin{center}
\small\renewcommand{\arraystretch}{1.4}
\begin{tabular}{l|l|l}
    \hline
    \rconcept{Тип}&\rconcept{Одноэлектронные числа}&\rconcept{Полный орбитальный момент}\\
    \hline
    E1&\(\Delta\ell=\pm1,\ \Delta n\) любое&
        \(\Delta L=0,\pm1;\quad0\not\leftrightarrow0\)\\
    M1&\(\Delta\ell=\Delta n=0\)&\(\Delta L=0\)\\
    E2&\(\Delta\ell=0,\pm2,\ \Delta n\) любое&
        \(\Delta L=0,\pm1,\pm2;\quad0\not\leftrightarrow0,1\)\\
    M2&\(\Delta\ell=\pm1\)&
        \(\Delta L=0,\pm1,\pm2;\quad0\not\leftrightarrow0,1\)\\
    \hline
\end{tabular}
\end{center}
Для E2 дополнительно \(\ell=0\not\leftrightarrow0\).
Правило \(\Delta L=0\) для многоэлектронного атома совместимо
с \(\Delta\ell=\pm1\) для отдельного электрона.

В чистой LS-связи:
\begin{equation*}
    \mathrm{E1,E2}:\ \Delta S=0;\qquad
    \mathrm{M1}:\ \Delta S=\Delta L=0,\quad
    \Delta J=\pm1\ \text{для линий внутри терма}.
\end{equation*}
При спин-орбитальном смешивании возникают слабые спин-меняющие линии.
Для ведущих каналов с \(\Delta S=\pm1\):
\begin{equation*}
    \begin{array}{c|c|c}
        \text{тип}&\Delta L&\text{исключение}\\ \hline
        \mathrm{E1,M1}&0,\pm1,\pm2&---\\
        \mathrm{E2}&0,\pm1,\pm2,\pm3&0\not\leftrightarrow0\\
        \mathrm{M2}&0,\pm1&0\not\leftrightarrow0
    \end{array}
\end{equation*}
Для E1, M1, E2 эти условия зависят от состава примесей.
Спиновая часть M2 допускает \(\Delta S=\pm1\) и между чистыми термами.
\rnote{Строгие ограничения по \(J\) и чётности сохраняются.}

\section{Спектры углерода и гелия}

\subsection{Спектр углерода}

Замкнутый остов \(1s^2 2s^2\) не влияет на термы. Для внешних электронов:
\begin{align*}
    2p^2:&\quad{}^3P_{0,1,2},\ {}^1D_2,\ {}^1S_0,
       &\pi&=(-1)^{1+1}=+1,\\
    2p3s:&\quad{}^3P^{\mathrm o}_{0,1,2},\ {}^1P^{\mathrm o}_1,
       &\pi&=(-1)^{1+0}=-1,\\
    2p3d:&\quad{}^{1,3}P^{\mathrm o},\ {}^{1,3}D^{\mathrm o},\ {}^{1,3}F^{\mathrm o},
       &\pi&=(-1)^{1+2}=-1.
\end{align*}
Индекс \(\mathrm o\) обозначает нечётность.
Для неэквивалентных электронов возможны \(S=0,1\);
\(p+s\) даёт \(L=1\), а \(p+d\) --- \(L=1,2,3\).

\begin{rbox}[left=10pt,right=10pt,top=10pt,bottom=8pt]{[]}
\centering
\captionsetup{hypcap=false}
\begin{tikzpicture}[x=1cm,y=1cm,>=Stealth,line width=0.7pt,
    every node/.style={font=\small}]
    \node at (1.4,3.7) {\(2p^2\), чётные};
    \node at (7.5,3.7) {\(2p3s\), нечётные};
    \foreach \y/\j in {0/0,0.5/1,1/2} {
        \draw (0,\y)--(2.2,\y) node[left,pos=0] {\({}^3P_{\j}\)};
    }
    \draw (0,2)--(2.2,2) node[left,pos=0] {\({}^1D_2\)};
    \draw (0,2.8)--(2.2,2.8) node[left,pos=0] {\({}^1S_0\)};
    \foreach \y/\j in {1.1/0,1.6/1,2.1/2} {
        \draw (6,\y)--(8.2,\y) node[right] {\({}^3P^{\mathrm o}_{\j}\)};
    }
    \draw (6,3.1)--(8.2,3.1) node[right] {\({}^1P^{\mathrm o}_1\)};
    \foreach \ya/\yb in {1.1/0.5,1.6/0,1.6/0.5,1.6/1,2.1/0.5,2.1/1} {
        \draw[->,rresultcolor] (6,\ya)--(2.2,\yb);
    }
    \draw[->,rconceptcolor] (6,3.1)--(2.2,2.8);
    \draw[->,rconceptcolor] (6,3.1)--(2.2,2);
\end{tikzpicture}
\captionof{figure}{E1-переходы \(2p3s\to2p^2\) в LS-приближении.
Голубые стрелки связывают триплеты, розовые --- синглеты.
Высоты уровней условны.}
\end{rbox}

Для \(2p3d\to2p^2\) E1-каналы группируются по термам:
\begin{equation*}
    \begin{array}{rcl}
        {}^3P^{\mathrm o},{}^3D^{\mathrm o}&\longrightarrow&{}^3P,\\
        {}^1P^{\mathrm o},{}^1D^{\mathrm o},{}^1F^{\mathrm o}
            &\longrightarrow&{}^1D,\\
        {}^1P^{\mathrm o}&\longrightarrow&{}^1S.
    \end{array}
\end{equation*}
Каждая компонента дополнительно удовлетворяет правилу по \(J\):
\({}^3P^{\mathrm o}_0\to{}^3P_1\) разрешён,
\({}^3P^{\mathrm o}_0\to{}^3P_0\) запрещён.

\subsection{Интеркомбинационные линии}

Спин-орбитальное взаимодействие смешивает состояния с одинаковыми
\(J\) и чётностью. Для преимущественно триплетного уровня гелия:
\begin{align*}
    |\widetilde{{}^3P^{\mathrm o}_1}\rangle
       &=\sqrt{1-|\eta|^2}\,|{}^3P^{\mathrm o}_1\rangle
         +\eta|{}^1P^{\mathrm o}_1\rangle,\qquad |\eta|\ll1,\\
    \langle{}^1S_0|\hat{\mathbf d}|\widetilde{{}^3P^{\mathrm o}_1}\rangle
       &=\eta\langle{}^1S_0|\hat{\mathbf d}|{}^1P^{\mathrm o}_1\rangle,\\
    w_{\text{интеркомб}}&\propto|\eta|^2|d_{\text{разр}}|^2.
\end{align*}
\(\eta\) --- коэффициент синглетной примеси.
\rconcept{Интеркомбинационная линия} связывает уровни разной
преобладающей мультиплетности; её малая интенсивность обусловлена
малостью примеси.

\clearpage
\subsection{Спектр гелия и метастабильные состояния}

\begin{rbox}[left=8pt,right=8pt,top=10pt,bottom=8pt]{[]}
\centering\captionsetup{hypcap=false}
\begin{tikzpicture}[x=0.92cm,y=0.80cm,>=Stealth,line width=0.7pt,
    every node/.style={font=\scriptsize}]
    \node at (3,5.9) {Синглеты: \(S=0\)};
    \node at (9.8,5.9) {Триплеты: \(S=1\)};
    \draw (0.6,0)--(12,0);
    \node[below] at (6,0) {\(1s^2\,{}^1S_0\)};
    \draw (0.6,2.9)--(2.5,2.9);
    \node[above] at (1.55,2.9) {\(1s2s\,{}^1S_0\)};
    \draw (3.6,3.8)--(5.5,3.8);
    \node[above] at (4.5,3.8) {\(1s2p\,{}^1P^{\mathrm o}_1\)};
    \draw (3.6,5.1)--(5.5,5.1);
    \node[above] at (4.5,5.1) {\(1s3d\,{}^1D_2\)};
    \draw (8.6,2)--(10.4,2);
    \node[below] at (9.25,2) {\(1s2s\,{}^3S_1\)};
    \foreach \y/\j in {4.7/0,4.3/1,3.9/2} {
        \draw (8.6,\y)--(10.4,\y) node[right] {\(J=\j\)};
    }
    \node[above] at (9.5,4.9) {\(1s2p\,{}^3P^{\mathrm o}_J\)};
    \draw[->,rresultcolor] (4.15,5.1)--(4.15,3.8)
        node[midway,left] {E1};
    \draw[->,rresultcolor] (3.6,3.8)--(2.5,2.9)
        node[midway,above left] {E1};
    \draw[->,rresultcolor] (4.55,3.8)--(4.55,0)
        node[pos=0.65,left] {E1};
    \foreach \x/\y in {8.8/4.7,9.35/4.3,9.9/3.9} {
        \draw[->,rresultcolor] (\x,\y)--(\x,2);
    }
    \node[rresultcolor] at (9.35,2.65) {E1};
    \draw[->,rconceptcolor] (1.2,2.9)--(1.2,0)
        node[midway,left] {E1+E1};
    \draw[->,rconceptcolor] (5.5,5.1)--(6.8,0)
        node[pos=0.55,right] {E2};
    \draw[->,rconceptcolor] (10.1,2)--(10.1,0)
        node[midway,left] {M1};
    \draw[->,rconceptcolor] (10.4,3.9)--(11.7,0)
        node[pos=0.6,right] {M2};
    \draw[->,rconceptcolor,dashed] (8.6,4.3)--(7.6,0)
        node[pos=0.7,left] {\(|\eta|^2\mathrm{E1}\)};
    \draw[->,rconceptcolor,dashed] (8.6,4.3)--(2.5,2.9);
    \node[rconceptcolor] at (6.85,3.65) {\(|\eta|^2\mathrm{E1}\)};
    \draw[<->,rconceptcolor,densely dotted]
        (5.5,3.8) to[bend left=45] (8.6,4.3);
    \node[rconceptcolor] at (6.8,4.95) {\(\hat H_{\rm so}\)};
\end{tikzpicture}
\captionof{figure}{Гелий: голубые стрелки --- разрешённые E1-каналы,
розовые --- слабые; пунктир --- смешивание. Масштаб условный.}
\end{rbox}

Характерные скорости разрешённых переходов:
\begin{equation*}
    \begin{array}{rcl}
    1s3d\,{}^1D_2\to1s2p\,{}^1P^{\mathrm o}_1
        &:&6.4\cdot10^7\ \text{с}^{-1},\\
    1s2p\,{}^1P^{\mathrm o}_1\to1s2s\,{}^1S_0
        &:&2\cdot10^6\ \text{с}^{-1},\\
    1s2p\,{}^3P^{\mathrm o}_J\to1s2s\,{}^3S_1
        &:&\sim10^7\ \text{с}^{-1}.
    \end{array}
\end{equation*}
При смешивании синглетного и триплетного \(P_1\)-уровней:
\begin{align*}
    w(\widetilde{{}^1P^{\mathrm o}_1}\to1s^2\,{}^1S_0)
        &\simeq1.8\cdot10^9(1-|\eta|^2)\ \text{с}^{-1},\\
    w(\widetilde{{}^3P^{\mathrm o}_1}\to1s^2\,{}^1S_0)
        &\simeq1.8\cdot10^9|\eta|^2\ \text{с}^{-1},\\
    w(\widetilde{{}^3P^{\mathrm o}_1}\to1s2s\,{}^1S_0)
        &\simeq3.0\cdot10^{-2}\ \text{с}^{-1}.
\end{align*}
Для триплетной группы \(1s2p\):
\(\Delta\widetilde\nu_{01}=0.99\ \text{см}^{-1}\),
\(\Delta\widetilde\nu_{12}=0.076\ \text{см}^{-1}\),
где \(\Delta E=hc\,\Delta\widetilde\nu\).

\rconcept{Метастабильный уровень} имеет большое время жизни
из-за отсутствия быстрых каналов распада.
Слабые переходы в основное состояние \(1s^2\,{}^1S_0\):
\begin{center}
\small\renewcommand{\arraystretch}{1.3}
\begin{tabular}{c|c|c}
    \hline
    \rconcept{Исходный уровень}&\rconcept{Канал}&\rconcept{\(w,\ \text{с}^{-1}\)}\\
    \hline
    \(1s2s\,{}^1S_0\)&E1+E1&51\\
    \(1s3d\,{}^1D_2\)&E2&\(\sim10^2\)\\
    \(1s2s\,{}^3S_1\)&M1&\(1.3\cdot10^{-4}\)\\
    \(1s2p\,{}^3P^{\mathrm o}_2\)&M2&0.33\\
    \hline
\end{tabular}
\end{center}
Слабый M1-распад \({}^3S_1\to{}^1S_0\) обусловлен поправками
к ведущему нерелятивистскому M1-оператору.

Для \(1s2s\,{}^1S_0\to1s^2\,{}^1S_0\):
\begin{equation*}
    \pi_i=\pi_f,\quad J_i=J_f=0
    \ \Rightarrow\ \mathrm{E1}\ \text{запрещён},\qquad
    \mathrm{E1+E1}\ \text{разрешён}.
\end{equation*}
Два фотона могут унести нулевой суммарный момент:
\begin{equation*}
    \hbar\omega_1+\hbar\omega_2=E_i-E_f,\qquad
    \tau\simeq\frac1{51}\approx0.020\ \text{с}.
\end{equation*}

\section{Масштабы взаимодействий и нарушение LS-связи}

\subsection{Спектр ртути}

С ростом заряда ядра усиливается спин-орбитальное взаимодействие.
Водородоподобная оценка в СГС:
\begin{equation*}
    mv^2\sim Z^2Ry\sim\frac{mZ^2e^4}{\hbar^2}
    \quad\Longrightarrow\quad
    \frac vc\sim Z\alpha,\qquad
    \alpha=\frac{e^2}{\hbar c}\simeq\frac1{137}.
\end{equation*}
Для внутренних электронов Hg \(Z\alpha\approx80/137\approx0.58\);
для внешних электронов существенно экранирование.

\begin{rbox}[left=8pt,right=8pt,top=10pt,bottom=8pt]{[]}
\centering\captionsetup{hypcap=false}
\begin{tikzpicture}[x=0.93cm,y=0.76cm,>=Stealth,line width=0.65pt,
    every node/.style={font=\scriptsize}]
    \draw[->] (-0.55,0)--(-0.55,7.8) node[above] {\(E\)};
    \foreach \x/\term in {0.6/{}^1S,2.2/{}^1P,3.8/{}^1D,5.4/{}^1F,
        7/{}^3S,8.6/{}^3P,10.2/{}^3D,11.8/{}^3F} {
        \node at (\x,8) {\(\term\)};
        \foreach \y in {7.05,7.25,7.4,7.52} {
            \draw (\x-0.42,\y)--(\x+0.42,\y);
        }
    }
    \draw (0.15,0)--(1.05,0) node[right] {\(6s^2\)};
    \foreach \x/\y/\orb in {0.6/4.7/7s,0.6/6.2/8s,2.2/3.1/6p,
        2.2/5.55/7p,2.2/6.35/8p,3.8/5.55/6d,3.8/6.5/7d,
        5.4/6.3/6f,7/4.25/7s,7/6/8s,8.6/5.15/7p,8.6/6.3/8p,
        10.2/5.45/6d,10.2/6.45/7d,11.8/6.3/6f} {
        \draw (\x-0.42,\y)--(\x+0.42,\y) node[right] {\(\orb\)};
    }
    \foreach \y/\j in {1.6/0,2.05/1,2.6/2} {
        \draw (8.18,\y)--(9.02,\y) node[right] {\({}^3P^{\mathrm o}_{\j}\)};
    }
    \node[left] at (8.15,2.05) {\(6p\)};
    \foreach \x/\y in {0.6/4.7,0.6/6.2,3.8/5.55,3.8/6.5} {
        \draw[->,rresultcolor] (\x,\y)--(2.2,3.1);
    }
    \draw[->,rresultcolor] (2.2,3.1)--(0.6,0);
    \draw[->,rresultcolor] (2.2,5.55)--(0.7,0);
    \draw[->,rresultcolor] (5.4,6.3)--(3.8,5.55);
    \draw[->,rresultcolor] (11.8,6.3)--(10.2,5.45);
    \draw[->,rresultcolor] (8.6,5.15)--(7,4.25);
    \foreach \y in {1.6,2.05,2.6} {
        \draw[->,rresultcolor] (7,4.25)--(8.6,\y);
        \draw[->,rresultcolor] (10.2,5.45)--(8.6,\y);
    }
    \draw[->,rconceptcolor,dashed] (7,6)--(2.2,3.1);
    \draw[->,rconceptcolor,dashed] (8.6,2.05)--(0.9,0);
\end{tikzpicture}
\captionof{figure}{Схема серий Hg. У возбуждённых уровней указан внешний
электрон; второй остаётся в \(6s\). Розовым пунктиром выделены
интеркомбинационные линии. Положения уровней условны.}
\end{rbox}

Примесь \({}^1P^{\mathrm o}_1\) открывает переход
\(6s6p\,{}^3P^{\mathrm o}_1\to6s^2\,{}^1S_0\).
Переходы \({}^3P^{\mathrm o}_{0,2}\to{}^1S_0\) остаются запрещёнными
как E1 по правилу для \(J\).

\clearpage
\subsection{Электростатическое и спин-орбитальное расщепление}

\begin{equation*}
    3d^3
    \ \xrightarrow{\text{электростатическое взаимодействие}}\
    {}^4F,\ {}^4P
    \ \xrightarrow{\text{спин-орбитальное взаимодействие}}\
    {}^4F_J,\ {}^4P_J.
\end{equation*}

\begin{rbox}[left=8pt,right=8pt,top=10pt,bottom=8pt]{[]}
\centering\captionsetup{hypcap=false}
\begin{tikzpicture}[x=1cm,y=0.82cm,>=Stealth,line width=0.7pt,
    every node/.style={font=\scriptsize}]
    \node at (5.6,7.65) {\(\mathrm{Cr}^{3+}:\ 3d^3\)};
    \draw[->] (-1.1,0)--(-1.1,3.9);
    \draw[->] (-1.1,4.4)--(-1.1,7.4) node[above] {\(E/(hc),\ \text{см}^{-1}\)};
    \draw (-1.25,4.08)--(-0.95,3.96)--(-1.25,3.84);
    \foreach \y/\j/\g/\energy in {0/3/4/0,0.95/5/6/244,2.15/7/8/561,3.6/9/10/956} {
        \draw (0,\y)--(3.2,\y);
        \node[above] at (0.35,\y) {\({}^4F_{\j/2}\)};
        \node[right] at (3.2,\y) {\(g=\g\)};
        \node[left] at (-1.2,\y) {\(\energy\)};
    }
    \foreach \y/\j/\g/\energy in {4.9/1/2/14072,5.8/3/4/14215,7/5/6/14481} {
        \draw (0,\y)--(3.2,\y);
        \node[above] at (0.35,\y) {\({}^4P_{\j/2}\)};
        \node[right] at (3.2,\y) {\(g=\g\)};
        \node[left] at (-1.2,\y) {\(\energy\)};
    }
    \foreach \a/\b/\v in {0/0.95/244,0.95/2.15/317,2.15/3.6/395,4.9/5.8/143,5.8/7/266} {
        \draw[<->,rresultcolor] (1.6,\a)--(1.6,\b)
            node[midway,right] {\(\v\)};
    }
    \draw (6,1.6)--(9.2,1.6) node[above,midway] {\(\overline E({}^4F)\)};
    \draw (6,6.1)--(9.2,6.1) node[above,midway] {\(\overline E({}^4P)\)};
    \draw[<->,rconceptcolor] (7.6,1.6)--(7.6,6.1)
        node[midway,right,align=left] {\(13774\ \text{см}^{-1}\)\\между термами};
\end{tikzpicture}
\captionof{figure}{Расщепление термов Cr\(^{3+}\). Разрыв оси отделяет
тонкую структуру от значительно большего межтермового интервала.}
\end{rbox}

Статистический вес и центр тяжести терма:
\begin{equation*}
    g_J=2J+1,\qquad
    \overline E_{LS}=\frac{\sum_J(2J+1)E_J}{\sum_J(2J+1)}.
\end{equation*}
В LS-приближении:
\begin{align*}
    E_J&=\overline E_{LS}+\frac{A_{LS}}2
       [J(J+1)-L(L+1)-S(S+1)],\\
    E_J-E_{J-1}&=A_{LS}J.
\end{align*}
Последнее соотношение --- \rconcept{правило интервалов Ланде}.
Отклонения реальных интервалов от него характеризуют неточность LS-модели.

\chapter{Взаимодействие атома с излучением}
\rsubtitle{Вероятности переходов и квантование поля}

\section{Гамильтониан и вероятность перехода}

\begin{equation*}
    \hat H=\underbrace{\hat H_E+\hat H_R}_{\hat H_0}+\hat H_I.
\end{equation*}
\(\hat H_E\) описывает атом, \(\hat H_R\) --- свободное поле,
\(\hat H_I\) --- взаимодействие. Для полной системы:
\begin{equation*}
    i\hbar\frac{\partial\Phi}{\partial t}=\hat H\Phi,\qquad
    \Phi(t)=\hat U(t,t_0)\Phi(t_0),\qquad
    \hat U(t,t_0)=e^{-i\hat H(t-t_0)/\hbar}.
\end{equation*}
Экспоненциальная запись верна при \(\partial\hat H/\partial t=0\).
Начальное \(|i\rangle\) и конечное \(|f\rangle\) состояния
включают состояние атома и числа фотонов:
\begin{equation*}
    A_{fi}=\langle f|\hat U|i\rangle,\qquad
    P_{fi}=|A_{fi}|^2,\qquad
    w_{fi}=\frac{dP_{fi}}{dt},\qquad [w_{fi}]=\text{с}^{-1}.
\end{equation*}
При постоянных скоростях независимых каналов распада:
\begin{equation*}
    N_i(t)=N_i(0)e^{-t/\tau_i},\qquad
    \frac1{\tau_i}=\sum_f w_{fi}.
\end{equation*}

\subsection{Дипольное приближение}

Если поле почти однородно на размере атома \(a\),
\begin{equation*}
    ka\ll1,\qquad k=\frac\omega c=\frac{2\pi}{\lambda},
    \qquad \hat H_I=-\hat{\mathbf d}\cdot\hat{\mathbf E}.
\end{equation*}
Для характерного атомного перехода в СГС:
\begin{equation*}
    a\sim\frac{\hbar^2}{me^2},\qquad
    \mathcal E\sim\frac{me^4}{\hbar^2},\qquad
    k\sim\frac{\mathcal E}{\hbar c}
    \quad\Rightarrow\quad ka\sim\alpha\ll1.
\end{equation*}
Через амплитуду векторного потенциала \(A\):
\begin{equation*}
    E\sim kA,\qquad H_I\sim eaE\sim e(ka)A\sim\alpha eA.
\end{equation*}
Малость \(ka\) обеспечивает дипольное приближение;
применимость теории возмущений требует также достаточно слабого поля.

\section{Временная теория возмущений}

\subsection{Адиабатическое включение}

Для задания невозмущённого состояния в далёком прошлом:
\begin{equation*}
    \hat H_I(t)=\hat V e^{\varepsilon t},\qquad
    t_0\to-\infty,\quad\varepsilon>0,\quad
    \lim_{t\to-\infty}\hat H_I(t)=0.
\end{equation*}
Предел \(\varepsilon\to+0\) берут после интегрирования.

В \rconcept{представлении взаимодействия}:
\begin{equation*}
    |\Phi_I(t)\rangle=e^{i\hat H_0t/\hbar}|\Phi(t)\rangle,\qquad
    \hat V_I(t)=e^{i\hat H_0t/\hbar}\hat V e^{-i\hat H_0t/\hbar}
    e^{\varepsilon t}.
\end{equation*}
Интегральное уравнение и его последовательные приближения:
\begin{align*}
    \hat U_I(t,-\infty)
    &=1-\frac{i}{\hbar}\int_{-\infty}^{t}
      \hat V_I(t_1)\hat U_I(t_1,-\infty)\,dt_1,\\
    \hat U_I^{(1)}
    &=-\frac{i}{\hbar}\int_{-\infty}^{t}\hat V_I(t_1)\,dt_1,\\
    \hat U_I^{(2)}
    &=-\frac1{\hbar^2}\int_{-\infty}^{t}dt_1
       \int_{-\infty}^{t_1}dt_2\,\hat V_I(t_1)\hat V_I(t_2),\\
    \hat U_I&=1+\hat U_I^{(1)}+\hat U_I^{(2)}+\cdots.
\end{align*}
Ограничение \(t_2\le t_1\) сохраняет временной порядок взаимодействий.

\subsection{Первый порядок: золотое правило Ферми}

Пусть \(V_{fi}=\langle f|\hat V|i\rangle\),
\(\omega_{fi}=(E_f-E_i)/\hbar\). Тогда
\begin{align*}
    A_{fi}^{(1)}
       &=-\frac{iV_{fi}}{\hbar}
         \int_{-\infty}^{t}e^{(i\omega_{fi}+\varepsilon)t_1}\,dt_1
        =-\frac{iV_{fi}}{\hbar}
          \frac{e^{(i\omega_{fi}+\varepsilon)t}}{\varepsilon+i\omega_{fi}},\\
    \frac{d|A_{fi}^{(1)}|^2}{dt}
       &=\frac{2|V_{fi}|^2}{\hbar^2}
          \frac{\varepsilon e^{2\varepsilon t}}
               {\varepsilon^2+\omega_{fi}^2}.
\end{align*}
Используя \(\lim_{\varepsilon\to+0}\varepsilon/(x^2+\varepsilon^2)
=\pi\delta(x)\), получаем \rconcept{золотое правило Ферми}:
\begin{equation*}
    w_{fi}^{(1)}=\frac{2\pi}{\hbar}|V_{fi}|^2\delta(E_f-E_i).
\end{equation*}
Дельта-функция сохраняет полную энергию атома и поля.
Для испускания одного фотона:
\begin{equation*}
    E_f=E_i\quad\Longleftrightarrow\quad
    \mathcal E_i-\mathcal E_f=\hbar\omega,
\end{equation*}
где \(\mathcal E_i,\mathcal E_f\) --- энергии самого атома.
При плотности конечных состояний \(\rho(E_f)\):
\begin{equation*}
    w_i=\frac{2\pi}{\hbar}\int
    |V_{fi}|^2\rho(E_f)\delta(E_f-E_i)\,dE_f.
\end{equation*}

\subsection{Второй порядок}

Два действия возмущения соединяются через промежуточное состояние:
\begin{equation*}
    i\xrightarrow{\ V_{li}\ }l\xrightarrow{\ V_{fl}\ }f,\qquad
    \sum_l|l\rangle\langle l|=1.
\end{equation*}
В нерезонансном случае:
\begin{align*}
    w_{fi}^{(2)}
    &=\frac{2\pi}{\hbar}
      \left|\sum_l\frac{V_{fl}V_{li}}{E_i-E_l+i0}\right|^2
      \delta(E_f-E_i)\\
    &=\frac{2\pi}{\hbar^3}
      \left|\sum_l\frac{V_{fl}V_{li}}{\omega_i-\omega_l+i0}\right|^2
      \delta(E_f-E_i),\qquad E_n=\hbar\omega_n.
\end{align*}
Складываются \emph{амплитуды} путей, поэтому возникает интерференция.
Промежуточные состояния виртуальные; на отдельном шаге энергия
может не сохраняться. При резонансе учитывают ширины уровней.
Второй порядок амплитуды даёт четвёртый порядок вероятности по \(\hat V\).

Для E1+E1-испускания возможны оба порядка рождения фотонов:
\begin{equation*}
    \mathcal A\propto\sum_n\left[
      \frac{(\mathbf e_2\cdot\mathbf d_{fn})(\mathbf e_1\cdot\mathbf d_{ni})}
           {\mathcal E_i-\mathcal E_n-\hbar\omega_1}
      +\frac{(\mathbf e_1\cdot\mathbf d_{fn})(\mathbf e_2\cdot\mathbf d_{ni})}
           {\mathcal E_i-\mathcal E_n-\hbar\omega_2}
    \right].
\end{equation*}
\(\mathbf e_1,\mathbf e_2\) --- поляризации;
\(\hbar(\omega_1+\omega_2)=\mathcal E_i-\mathcal E_f\).

\section{Операторы переходов и квантование поля}

\subsection{Атом и фотонные моды}

Оператор перехода атома из \(j\) в \(i\):
\begin{equation*}
    \hat b_i^\dagger\hat b_j=|i\rangle\langle j|,\qquad
    |i\rangle\langle j|k\rangle=\delta_{jk}|i\rangle.
\end{equation*}
Равенство понимается в одноатомном секторе. Разложение диполя:
\begin{equation*}
    \hat{\mathbf d}=\sum_{i,j}\mathbf d_{ij}\hat b_i^\dagger\hat b_j,
    \qquad \mathbf d_{ij}=\langle i|\hat{\mathbf d}|j\rangle.
\end{equation*}
Для моды \(k\) с \(N_k\) фотонами:
\begin{equation*}
    \hat a_k|N_k\rangle=\sqrt{N_k}|N_k-1\rangle,\qquad
    \hat a_k^\dagger|N_k\rangle=\sqrt{N_k+1}|N_k+1\rangle.
\end{equation*}
\(a_k\) уничтожает фотон, \(a_k^\dagger\) рождает.
При вещественных модовых коэффициентах в точке атома:
\begin{align*}
    \hat{\mathbf E}(t)
      &=i\sum_k\mathcal E_k\mathbf e_k
        \bigl(\hat a_ke^{-i\omega_kt}-\hat a_k^\dagger e^{i\omega_kt}\bigr),\\
    \hat H_I&=-\hat{\mathbf d}\cdot\hat{\mathbf E}.
\end{align*}
\(\mathbf e_k\) --- поляризация, \(\mathcal E_k\) --- амплитудный
коэффициент моды; пространственная фаза включена в её определение.

\subsection{Двухуровневый атом}

При \(E_2>E_1\) введём
\begin{equation*}
    \hat\sigma_+=|2\rangle\langle1|,\quad
    \hat\sigma_-=|1\rangle\langle2|,\quad
    \hat d=d_{12}(\hat\sigma_-+\hat\sigma_+).
\end{equation*}
Здесь \(d_{11}=d_{22}=0\) по чётности, фаза выбрана так, что
\(d_{12}\) вещественен. Для одной моды получаются четыре члена:
\begin{equation*}
    (\hat\sigma_-+\hat\sigma_+)(\hat a-\hat a^\dagger)
    =\underbrace{\hat\sigma_-\hat a}_{\text{а}}
     +\underbrace{\hat\sigma_+\hat a}_{\text{б}}
     -\underbrace{\hat\sigma_-\hat a^\dagger}_{\text{в}}
     -\underbrace{\hat\sigma_+\hat a^\dagger}_{\text{г}}.
\end{equation*}

\begin{rbox}[left=10pt,right=10pt,top=10pt,bottom=8pt]{[]}
\centering
\captionsetup{hypcap=false}
\begin{tikzpicture}[x=1cm,y=1cm,>=Stealth,line width=0.8pt,
    every node/.style={font=\small}]
    \foreach \xx/\yy/\lab/\ini/\fin in {0/2.3/а/2/1,6/2.3/б/1/2,0/0/в/2/1,6/0/г/1/2} {
        \begin{scope}[shift={(\xx,\yy)}]
            \node at (-0.5,0.4) {(\lab)};
            \draw[->] (0,0)--(1.5,0) node[midway,below] {\(\ini\)};
            \draw[->] (1.5,0)--(3.3,0) node[midway,below] {\(\fin\)};
            \fill (1.5,0) circle (1.8pt);
        \end{scope}
    }
    \foreach \xx in {0,6} {
        \begin{scope}[shift={(\xx,2.3)}]
            \draw[rconceptcolor] plot[domain=0:1,samples=50]
                ({0.3+1.2*\x},{1.1*(1-\x)+0.08*sin(1440*\x)});
            \draw[->,rconceptcolor] (1.35,0.14)--(1.5,0);
            \node[rconceptcolor] at (0.25,1.4) {\(k\)};
        \end{scope}
        \begin{scope}[shift={(\xx,0)}]
            \draw[rconceptcolor] plot[domain=0:1,samples=50]
                ({1.5+1.2*\x},{1.1*\x+0.08*sin(1440*\x)});
            \draw[->,rconceptcolor] (2.7,1.1)--(2.85,1.24);
            \node[rconceptcolor] at (3.0,1.4) {\(k\)};
        \end{scope}
    }
\end{tikzpicture}
\captionof{figure}{Диаграммы четырёх членов взаимодействия.
Прямая линия --- атом, волнистая --- фотон, вершина --- одно действие
взаимодействия. Время направлено слева направо.}
\end{rbox}

\begin{center}
\small\renewcommand{\arraystretch}{1.35}
\begin{tabular}{c|c|l}
    \hline
    &\rconcept{Изменение состояния}&\rconcept{Процесс}\\
    \hline
    а&\(|2,N\rangle\to|1,N-1\rangle\)&понижение атома и уничтожение фотона\\
    б&\(|1,N\rangle\to|2,N-1\rangle\)&поглощение\\
    в&\(|2,N\rangle\to|1,N+1\rangle\)&испускание\\
    г&\(|1,N\rangle\to|2,N+1\rangle\)&возбуждение атома и рождение фотона\\
    \hline
\end{tabular}
\end{center}
При резонансе \(\omega\simeq\omega_{21}=(E_2-E_1)/\hbar\)
реальные однофотонные процессы описывают б и в.
Противовращающиеся члены а и г осциллируют с частотой
\(\omega+\omega_{21}\); в слабом резонансном поле ими пренебрегают,
но они участвуют в виртуальных процессах.

Множители \(\sqrt{N_k}\) и \(\sqrt{N_k+1}\) дают
\begin{equation*}
    w_{\text{погл}}\propto N_k,\qquad
    w_{\text{исп}}\propto
    \underbrace{N_k}_{\text{вынужденное}}
    +\underbrace{1}_{\text{спонтанное}}.
\end{equation*}
Спонтанное испускание возможно и при \(N_k=0\).

\clearpage
\section{Гамильтониан через векторный потенциал}

В СГС, без явного спинового члена:
\begin{equation*}
    \hat H=\sum_{a=1}^{N_e}\frac1{2m}
      \left[\hat{\mathbf p}_a-\frac{q_e}{c}
      \hat{\mathbf A}(\mathbf r_a)\right]^2+U+\hat H_R.
\end{equation*}
\(U\) --- кулоновская энергия, \(N_e\) --- число электронов,
\(q_e=-e\). Раскрывая квадрат:
\begin{align*}
    \hat H_E&=\sum_a\frac{\hat{\mathbf p}_a^2}{2m}+U,\\
    \hat H_I&=-\frac{q_e}{2mc}\sum_a
      (\hat{\mathbf p}_a\cdot\hat{\mathbf A}_a+
       \hat{\mathbf A}_a\cdot\hat{\mathbf p}_a)
      +\frac{q_e^2}{2mc^2}\sum_a\hat{\mathbf A}_a^2,
      \qquad \hat{\mathbf A}_a=\hat{\mathbf A}(\mathbf r_a).
\end{align*}
В кулоновской калибровке \(\boldsymbol\nabla\cdot\mathbf A=0\):
\begin{equation*}
    \hat H_I=
      \underbrace{-\frac{q_e}{mc}\sum_a
        \hat{\mathbf A}_a\cdot\hat{\mathbf p}_a}_{\hat H_I^{(A)}}
      +\underbrace{\frac{q_e^2}{2mc^2}\sum_a
        \hat{\mathbf A}_a^2}_{\hat H_I^{(A^2)}}.
\end{equation*}
Структура фотонных операторов:
\begin{equation*}
    \begin{array}{ccl}
        \mathbf p\cdot\mathbf A&\sim a,\ a^\dagger&
            \text{однофотонное поглощение и испускание},\\
        \mathbf A^2&\sim aa,\ a^\dagger a^\dagger,\ a^\dagger a,\ aa^\dagger&
            \text{двухфотонные процессы и рассеяние}.
    \end{array}
\end{equation*}
Двухфотонные переходы возникают также во втором порядке по
\(\hat H_I^{(A)}\). В строгом дипольном приближении
\begin{equation*}
    \mathbf A(\mathbf r_a)\approx\mathbf A(0),\qquad
    \sum_a\mathbf A_a^2=N_e\mathbf A^2(0),\qquad
    \langle f_{\rm at}|\mathbf A^2(0)|i_{\rm at}\rangle=0
    \quad(f_{\rm at}\ne i_{\rm at}).
\end{equation*}
Поэтому E1+E1-переход между разными атомными уровнями описывается
двумя линейными взаимодействиями через промежуточные состояния.

\end{document}
